\section{mamaretrieval: A Chunk-Level Relevance Benchmark}
\label{sec:mamaretrieval}

mamaretrieval scores how well a retriever finds the guideline passages that
answer a clinical query, at the chunk granularity a RAG system actually
retrieves. Building it means choosing what to retrieve over (the corpus), what
queries to ask, how to label a (query, chunk) pair's relevance, and how to keep
those labels trustworthy at scale. We start with the corpus, then take query
construction, the relevance rubric, pooling, and judge calibration in turn.

\subsection{The guideline corpus}
\label{sec:corpus}

mamaretrieval is built over the same maternal-health guideline corpus the
deployed system retrieves from, which the companion system
paper~\cite{ren2026mamai} describes in full---its sources, construction, and
packaging. We give here only what the benchmark depends on: the corpus's scale,
its source-quality tiers, and the version coupling that makes the benchmark and
its corpus one artifact.

\paragraph{Scale and sources.}
The corpus bundle (v0.2.0) holds 63{,}650 section-aware chunks drawn from 87
documents spanning international clinical guidance (WHO, NICE, MSF, the
Hesperian midwifery references, and FIGO/EmONC materials), national and local
protocols (the Tanzania and Zanzibar Ministry of Health guidelines), and
standard obstetric and midwifery reference texts. Each chunk carries a canonical
header recording its source, page, and a stable chunk identifier.

\paragraph{Source-quality tiers.}
Sources are rated into quality tiers (from \emph{very high} down to
\emph{low--moderate}) by their clinical depth and relevance to the
Zanzibar OBGYN and midwifery scope. Query construction (below) uses these tiers
as sampling weights, drawing proportionally more from the clinically richest
sources.

\paragraph{Version coupling.}
Chunk identifiers are content-derived---a hash of the chunk text---so they
survive source renames and page drift but change whenever a chunk's text
changes. mamaretrieval references chunks by these identifiers, which couples it
to a specific \emph{chunk set} rather than to a specific embedding. The
distinction is not hypothetical: the deployed system moved from corpus bundle
v0.2.0 to v0.3.0, but that update only re-embedded the corpus---swapping the
on-device embedder---and left all 63{,}650 chunks and their identifiers
untouched, so every mamaretrieval label applies to the deployed v0.3.0 corpus
unchanged. A genuine re-chunking would change the identifiers and require
re-running the labelling pipeline; we therefore treat the
(chunk set, queries, labels) triple as a single versioned artifact rather than
versioning queries and labels on their own.

\subsection{Query construction}
\label{sec:mr-queries}

A retrieval benchmark needs queries, and for this corpus there are none to
mine---no search logs, no question bank tied to these guidelines. We generate
them. Each query in mamaretrieval is written by an LLM from a single corpus
chunk, asking a clinical question that the chunk answers; the chunk that seeded a
query is recorded as a known-relevant anchor for it.

The queries come out of a three-step funnel over the corpus:
\begin{enumerate}
  \item \textbf{Tier-weighted sampling} draws 4{,}540 chunks, with each source's
  quota set by its quality tier (\S\ref{sec:corpus})---the clinically richest
  sources contribute proportionally more, so the query set inherits the corpus's
  quality signal rather than its raw size distribution.
  \item An \textbf{answerability gate} passes each sampled chunk to
  \texttt{Qwen3.6-27B-FP8} (reasoning on, temperature 0), prompted to keep the
  chunk only if a clinician could answer a specific question from it alone and,
  when so, to write exactly that question.\footnote{The full filter-and-generate
  prompt ships with the construction code (\S\ref{sec:access}).}
  \item The kept chunks become \textbf{query records} with stable identifiers.
\end{enumerate}
The gate keeps 3{,}185 of the 4{,}540 chunks (about 70\%), and the 30\% it
discards are themselves informative: assessment rubrics, learning objectives, and
table-of-contents fragments are rejected because no answerable clinical question
follows from them.

Two limits of this design are worth stating up front. First, only \emph{per
chunk} queries---one chunk, one question---are generated in this release; the
specification also allows synthesis queries spanning several chunks and
adversarial variants, but those are deferred (\S\ref{sec:limitations}). Second,
and more consequentially, a query written from a chunk tends to be phrased like
that chunk, which gives dense embedding retrievers an artificial edge over
lexical ones on the very chunk that seeded the query. We do not correct for this
anchor-text effect; we disclose it, and it shapes which metrics we trust
(\S\ref{sec:mr-judge}).

\subsection{The graded relevance rubric}
\label{sec:mr-rubric}

With queries and a corpus, the benchmark needs a relevance label for each
(query, chunk) pair: how useful is this chunk for answering this query? Our first
rubric made that judgement binary---relevant or not. It saturated. Take a query
asking for alternative iron and folic-acid dosages: a footnote giving the exact
substitute formulation, a dosing section padded with storage advice, and an
adherence-counselling note that never states a dose are, by a binary test,
equally relevant---each is on topic, meaningful, and broadly actionable. A
benchmark that cannot separate them cannot tell a retriever that surfaces the
dose from one that surfaces the filler.

So we replaced the binary label with a graded one. Each (query, chunk) pair is
scored
\[
  \mathrm{score} = D_1 \times (D_2 + D_3 + D_4) \in \{0,\dots,6\},
\]
one gate and three graded dimensions:
\begin{description}
  \item[$D_1$ --- topic (0/1).] Does the chunk address the query's clinical
  question in the same context---the same timing window (antenatal,
  intrapartum, postpartum)? A zero here forces a score of zero and halts
  scoring, which also keeps the judge from spending tokens on off-topic chunks.
  \item[$D_2$ --- meaningful content (0/1/2).] How much clinical depth the chunk
  carries, from a bare mention to a multi-concept explanation.
  \item[$D_3$ --- actionable guidance (0/1/2).] How specific the guidance is:
  none, a general instruction (``give a uterotonic''), or an exact one
  (``10 IU oxytocin IM'').
  \item[$D_4$ --- density (0/1/2).] What fraction of the chunk is useful for
  \emph{this} query, penalising a relevant passage buried in unrelated text.
\end{description}
A structural rule ties them together: actionable guidance presupposes some
meaningful content ($D_3 > 0 \Rightarrow D_2 > 0$).

The design follows established relevance theory rather than inventing one.
Decomposing relevance into orthogonal dimensions combined by a fixed rule is the
structure of the TREC Precision Medicine track~\cite{roberts2017precision}; the
split between topical, cognitive, and situational relevance is
Saracevic's~\cite{saracevic2007relevance}, with $D_1$ the topical gate and
$D_3$/$D_4$ the situational utility; and decomposing relevance this way has been
shown to raise inter-judge agreement over holistic scoring~\cite{farzi2025criteria}.
Two LLM-judge practices shape the rest: a clear topical gate, as in
UMBRELA~\cite{upadhyay2024umbrela}, and forcing the judge to reason before it
commits to a score, as in G-Eval~\cite{liu2023geval}---the structured-output
schema emits the reasoning first and the four dimension values last. The judge
is \texttt{Qwen3.5-397B-A17B-FP8} at temperature 0;
\S\ref{sec:mr-judge} reports how well its labels hold up.

\begin{table}[t]
\centering
\small
\setlength{\tabcolsep}{4pt}
\renewcommand{\arraystretch}{1.2}
\caption{The graded rubric on one query---``What are the alternative dosages for
iron and folic acid supplementation?''---as scored by the production judge over
four \emph{on-topic} chunks ($D_1{=}1$ throughout). Chunk text is verbatim;
sources in brackets (MSF, the MSF obstetric manual; WHO~PCPNC, the WHO
pregnancy/childbirth guide; WHO~ANC, the 2016 WHO antenatal-care guideline). A
binary label cannot separate the first three; the graded score ranks them 6, 4,
and 2.}
\label{tab:rubric-example}
\begin{tabular}{@{}p{0.49\columnwidth}@{\hspace{5pt}}ccccc@{}}
\toprule
Retrieved chunk (verbatim) & $D_1$ & $D_2$ & $D_3$ & $D_4$ & Score \\
\midrule
\emph{[MSF]}~``200 mg ferrous sulfate (65 mg elemental iron) + 400 micrograms
folic acid tablets may be replaced by 185 mg ferrous fumarate (60 mg elemental
iron) + 400 micrograms folic acid tablets.'' & 1 & 2 & 2 & 2 & \textbf{6} \\
\addlinespace
\emph{[WHO PCPNC]}~``Give iron and folic acid to all pregnant, postpartum and
post-abortion women: routinely once daily in pregnancy and until 3 months after
delivery or abortion; twice daily for anaemia (double dose). Check the supply at
each visit and dispense 3 months' supply; advise to store iron safely.'' & 1 & 2
& 1 & 1 & \textbf{4} \\
\addlinespace
\emph{[WHO PCPNC]}~``Motivate on adherence with treatments. Explore local
perceptions about iron treatment (e.g.\ that making more blood will worsen
bleeding, or that iron will make the baby too large).'' & 1 & 1 & 1 & 0 &
\textbf{2} \\
\addlinespace
\emph{[WHO ANC]}~``Resources. Intermittent iron and folic acid supplementation
might cost a little less than daily supplementation due to the lower total weekly
dose of iron.'' & 1 & 0 & 0 & 0 & \textbf{0} \\
\bottomrule
\end{tabular}
\end{table}

Table~\ref{tab:rubric-example} shows the graded rubric on that query. All four
chunks pass the topic gate; the first three also clear a binary rubric's bar
identically---each is meaningful and actionable---yet the graded rubric ranks
them apart. The footnote giving the exact alternative formulation scores 6,
specific and entirely on-query; a dosing section half-taken up by storage and
supply logistics scores 4 ($D_4=1$); and an adherence note that never states a
dose scores 2. The fourth chunk is on topic but pure cost commentary, and drops
to 0---the rubric counts resource text as no meaningful clinical content
($D_2=0$). The judge's reasoning ships with every label and makes such calls
explicit: of that last chunk, it notes the passage sits ``literally under a
\dots\ Resources heading and discusses cost,'' and scores $D_2=0$.

\subsection{Pooling and label completeness}
\label{sec:mr-pooling}

Scoring every chunk against every query is infeasible---63{,}650 chunks times
3{,}185 queries---so mamaretrieval labels a \emph{pool}: for each query, the
union of the top-ranked chunks returned by several retrievers, with every pooled
pair judged and every unpooled chunk treated as not-relevant. This is the
standard information-retrieval pooling protocol, and its one risk is well known:
if the pool misses relevant chunks, the labels under-count them. We measured that
risk rather than assume it away.

The first pool took the top-10 chunks from three retrievers (a lexical baseline,
a biomedical encoder, and a strong general-purpose embedder)---about 24.7
candidates per query and 78{,}571 judged pairs. A completeness audit on a
100-query sample, pooling a wider panel of six retrievers, found that this first
pool captured only about half of the relevant material (lenient recall near
0.49): half the relevant chunks were being labelled not-relevant by omission.

That finding drove the released label set. mamaretrieval v0.2.0 pools the top-20
chunks from all six retrievers across all 3{,}185 queries, yielding 230{,}964
graded labels---the pool the companion system paper draws its label provenance
from. A residual gap remains: chunks that no retriever ranks in its top-20 are
still unjudged, so even this pool is near-saturated at the top of the ranking
rather than exhaustive, and absolute coverage-sensitive numbers should be read
with that in mind (\S\ref{sec:mr-judge}).

\subsection{Judge calibration and metric guidance}
\label{sec:mr-judge}

A single LLM produces every relevance label, so the benchmark is only as
trustworthy as that judge. We calibrate it against a frontier reference model
(Claude~Opus~4.7) on a 62-pair pilot: the production judge agrees on 95\% of
pairs at the lenient cut (score $\ge 3$, any useful content) and 85\% at the
strict cut (score $\ge 5$, complete and specific), and a stratified manual
spot-check of its reasoning traces found its calls defensible, including on hard
cases such as drug-name collisions across clinical contexts. Those two
cuts---lenient $\ge 3$ and strict $\ge 5$---are the relevance thresholds the
benchmark reports against.

We are deliberate about what this establishes. It is an agreement check against
another model, not against clinicians: it shows the judge is a reasonable
relevance classifier and that retriever \emph{rankings} are stable under it, but
it does not make any single label ground truth. Each label is one calibrated
opinion, and absolute, per-row numbers carry more uncertainty than relative
comparisons.

That uncertainty, together with the pool incompleteness of
\S\ref{sec:mr-pooling}, is why we guide consumers on metric choice. Metrics that
ask only whether a relevant chunk is present---Hit~Rate@$k$ and MRR---are robust
to a pool that misses some relevant chunks. nDCG@$k$ is moderately sensitive,
since its discounting assumes the labels are complete. Recall@$k$ and
Precision@$k$ are the most sensitive, because the missing-label problem changes
their denominators directly. We recommend leading with the robust metrics and
reporting the rest with the completeness caveat attached---the order the
companion system paper follows when it evaluates retrievers on mamaretrieval.
